% ICCSDI 2026 conference manuscript — ACIS-X
%
% International Conference on Computational Science and Data Intelligence
% (ICCSDI 2026), 11--12 December 2026, NMIMS Mumbai, India (Hybrid Mode)

\documentclass[pdflatex,sn-mathphys-num]{sn-jnl}

%%%% Packages used by the ICCSDI 2026 conference template
\usepackage{graphicx}
\usepackage{multirow}
\usepackage{amsmath,amssymb,amsfonts}
\usepackage{amsthm}
\usepackage{mathrsfs}
\usepackage[title]{appendix}
\usepackage{xcolor}
\usepackage{textcomp}
\usepackage{manyfoot}
\usepackage{booktabs}
\usepackage{algorithm}
\usepackage{algorithmicx}
\usepackage{algpseudocode}
\usepackage{listings}
\usepackage{array}

%%%% Space-efficient float parameters
\setlength{\textfloatsep}{6pt plus 2pt minus 2pt}
\setlength{\floatsep}{6pt plus 2pt minus 2pt}
\setlength{\intextsep}{6pt plus 2pt minus 2pt}
\setlength{\abovecaptionskip}{4pt}
\setlength{\belowcaptionskip}{2pt}
\setlength{\abovedisplayskip}{5pt plus 2pt minus 2pt}
\setlength{\belowdisplayskip}{5pt plus 2pt minus 2pt}
\setlength{\abovedisplayshortskip}{2pt plus 1pt minus 1pt}
\setlength{\belowdisplayshortskip}{2pt plus 1pt minus 1pt}

%%%% Optional theorem environments
\theoremstyle{thmstyleone}
\newtheorem{theorem}{Theorem}
\newtheorem{proposition}[theorem]{Proposition}

\theoremstyle{thmstyletwo}
\newtheorem{example}{Example}
\newtheorem{remark}{Remark}

\theoremstyle{thmstylethree}
\newtheorem{definition}{Definition}

\raggedbottom

\begin{document}

%%%% Manuscript title
\title[ACIS-X: Multi-Agent Credit Intelligence]{Autonomous Credit Intelligence: An Event-Driven, Self-Healing Multi-Agent Framework for Real-Time B2B Risk Assessment}

%%%% Authors and affiliations
\author*[1]{\fnm{Nikhil} \sur{Nagendra}}\email{nikhilnag98@gmail.com}

\author[2]{\fnm{Pallavi} \sur{K V}}\email{pallavi.kv@dsat.edu.in}

\author[1]{\fnm{Bhoomi} \sur{Joshi}}

\author[1]{\fnm{Johan} \sur{K George}}

\author[1]{\fnm{Nitin Bharadwaj} \sur{MVS}}

\author[1]{\fnm{Prasanna B} \sur{Acharya}}

\affil*[1]{\orgdiv{Department of Computer Science and Engineering},
  \orgname{AMC Engineering College},
  \orgaddress{\city{Bengaluru}, \postcode{560083}, \country{India}}}

\affil[2]{\orgdiv{Department of Information Science and Engineering},
  \orgname{Dayanand Sagar Academy of Technology and Management},
  \orgaddress{\city{Bengaluru}, \postcode{560082}, \country{India}}}

%%%% Abstract
\abstract{Intraday B2B credit monitoring requires continuous synthesis of payment behaviour, balance-sheet fundamentals, and litigation signals. Given the proprietary confidentiality of commercial enterprise trade ledgers, evaluating streaming credit architectures necessitates reproducible macroeconomic stress simulation. This paper presents ACIS-X, a reactive event-driven multi-agent framework on Apache Kafka that continuously maintains a live per-customer risk score. Specialized agents coordinate over partitioned Kafka topics across five core layers, an autonomous model governance plane, and a self-healing plane. ACIS-X combines Kafka Exactly-Once Semantics (EOS) with idempotent database sinks, three-tier signal fusion, SHAP audit logging, streaming drift tracking via Population Stability Index (PSI), and heartbeat-driven recovery from committed offsets. Evaluated under a jump-diffusion stress harness, ACIS-X demonstrates sub-50\,ms streaming latency (P95 16.2\,ms), near-linear consumer scaling, and autonomous fault recovery with storm prevention.}

%%%% Keywords
\keywords{Multi-agent systems, Event-driven architecture, Credit risk intelligence, Apache Kafka, Exactly-Once Semantics, Self-healing systems}

\maketitle

\begin{center}
\small
\textbf{International Conference on Computational Science and Data Intelligence (ICCSDI 2026)}\\
11--12 December 2026, NMIMS Mumbai, India $\vert$ Hybrid Mode
\end{center}
\vspace{0.2em}

\renewcommand{\theHtable}{\thesection.\arabic{table}}

%%% ====================================================================
\section{Introduction}\label{sec:introduction}
%%% ====================================================================

Enterprise credit risk monitoring demands continuous evaluation of payment dynamics, financial health, and litigation. Event streaming platforms like Apache Kafka~\cite{kreps2011kafka} enable real-time risk updates, but recent LLM-based multi-agent systems~\cite{jajoo2025masca} introduce multi-second latencies ($>2$\,s) and non-deterministic outputs unsuitable for sub-50\,ms financial streaming. Conventional Kafka pipelines perform static transaction scoring without agentic autonomy or drift governance.

A fundamental challenge is that real B2B invoice ledgers and defaults are confidential commercial secrets. We therefore separate \textit{architectural validation} (the core contribution) from \textit{empirical default forecasting}. We introduce ACIS-X (Autonomous Credit Intelligence System, Extended), an event-driven multi-agent framework grounded in the Actor model, with the following contributions:
\begin{itemize}
\setlength{\itemsep}{1pt}\setlength{\parskip}{0pt}
\item An event-driven multi-agent architecture on Kafka with autonomous self-healing and model governance.
\item Kafka stream-level EOS combined with idempotent persistence and action deduplication for external sinks.
\item Three-tier signal fusion integrating behavioural metrics, financial ratios, and litigation with deterministic risk policy guardrails.
\item Streaming model governance via PSI drift tracking and online challenger benchmarking, inspired by SR~11-7.
\item Systems benchmarks measuring throughput, scaling, P95/P99 latency, and fault recovery under a jump-diffusion stress harness.
\end{itemize}

%%% ====================================================================
\section{Related Work and Research Gap}\label{sec:related}
%%% ====================================================================

Corporate credit scoring---from Z-scores~\cite{altman1968financial} to explainable ML~\cite{chang2025explainable,nwafor2024hybrid}---relies on offline batch feature stores incapable of capturing intraday delinquency transitions. Kafka event sourcing is adopted for sub-second fraud scoring~\cite{kesarpu2025kafka,parnerkar2025eda}, yet existing architectures evaluate isolated transactions rather than stateful, multi-source credit profiles combining balance-sheet ratios and litigation.

Cognitive MAS frameworks---MASCA~\cite{jajoo2025masca}, Kubam's credit agent~\cite{kubam2026agentic}, Kou et al.~\cite{kou2025learning}---applied LLMs and MARL to credit assessment, but their multi-second execution times and prompt non-determinism are incompatible with high-throughput, low-latency streaming. Existing autonomic resilience~\cite{beereddy2025selfhealing,kephart2003autonomic} recovers at coarse container boundaries rather than providing per-agent heartbeat recovery with offset replay. Table~\ref{tab:comparison} positions ACIS-X at the intersection of streaming fusion, deterministic compliance, autonomous governance, and fine-grained agent recovery.

\begin{table}[htbp]
\caption{Architectural positioning across credit risk paradigms.}\label{tab:comparison}
\centering\footnotesize
\begin{tabular}{@{}lcccc@{}}
\toprule
\textbf{Dimension} & \textbf{Batch XAI~\cite{chang2025explainable}} & \textbf{Stream~\cite{parnerkar2025eda}} & \textbf{Cog.\ MAS~\cite{jajoo2025masca}} & \textbf{ACIS-X} \\
\midrule
Latency       & Hours--Days  & $<$1\,s     & $>$2\,s        & P95 16\,ms \\
Signal Scope  & Fin.\ ratios & Single txn  & Text/reports   & 3-tier fusion \\
Exec.\ Model  & Batch        & Stream proc.& LLM/MARL       & Reactive MAS \\
Atomicity     & DB ACID      & AL-once/EOS & None           & EOS + PK \\
Fault Recov.  & Cold restart & Node failover& Prompt retry  & HB + offset \\
Drift Gov.    & Offline      & Ext.\ dash. & Manual         & PSI + Challenger \\
\botrule
\end{tabular}
\end{table}

%%% ====================================================================
\section{System Architecture}\label{sec:architecture}
%%% ====================================================================

ACIS-X maintains a live risk score for each customer $c_i$:
\begin{equation}
\mathcal{R}(c_i,t) = \mathcal{F}\bigl(\mathcal{B}(c_i,t),\;\mathcal{E}(c_i,t),\;\mathcal{L}(c_i,t)\bigr) \in [0,1],
\label{eq:risk}
\end{equation}
where $\mathcal{B}$ is behavioural history, $\mathcal{E}$ external financial indicators, and $\mathcal{L}$ litigation exposure. Design targets include sub-second processing, usable scores under partial signals, zero silent loss of committed work, and idempotent duplicate handling.

The architecture comprises six layers on a common Kafka event bus: Layer~1 ingests customers, invoices, payments, and behavioural indicators; Layer~2 enriches via financial APIs and litigation scraping; Layer~3 performs payment risk prediction and composite scoring; Layer~4 enforces credit policy and collections; Layer~5 persists state; Layer~6 tracks drift and benchmarks challengers. A system plane manages time ticks, monitoring, self-healing, and dead-letter-queue isolation. All events are keyed by \texttt{customer\_id}, ensuring per-customer partition affinity and temporal ordering without distributed join coordinators.

Table~\ref{tab:agents} catalogues every agent, its layer, and operational role. Each operates as an isolated OS process with a dedicated consumer group.

\begin{table}[htbp]
\caption{ACIS-X agent inventory: layer and operational role.}\label{tab:agents}
\centering\footnotesize
\begin{tabular}{@{}llp{5.2cm}@{}}
\toprule
\textbf{Agent} & \textbf{L} & \textbf{Role} \\
\midrule
ScenarioGenerator  & 1 & Emits reproducible synthetic/jump-diffusion macro-stress streams. \\
CustomerState      & 1 & Rebuilds behavioural metrics (OTR, delay, utilisation). \\
ExternalData       & 2 & Enriches and caches financial ratios (D/E, ROE, PE). \\
LitigationScraping & 2 & Quantifies legal exposure from court/news feeds. \\
Aggregator         & 2 & Fuses Tier~2 financial and litigation signals ($0.6/0.4$). \\
PaymentPrediction  & 3 & RF scoring + guardrails + SHAP attributions. \\
RiskScoring        & 3 & Refines base score with behavioural \& temporal velocity. \\
OverdueDetection   & 4 & Time-tick driven delinquency detection. \\
CreditPolicy       & 4 & Maps risk tiers to operational intents. \\
Collections        & 4 & Idempotent collection actions; storm prevention. \\
DBAgent            & 5 & Persists events (SQLite/Postgres) for recovery. \\
Query / Memory     & 5 & In-memory reads with DB fallback. \\
ModelGovernance    & 6 & PSI/KS drift tracking; challenger benchmarking. \\
\botrule
\end{tabular}
\end{table}

%%% ====================================================================
\section{Methodology}\label{sec:methods}
%%% ====================================================================

\subsection{Three-Tier Signal Fusion}\label{subsec:fusion}

\textbf{Tier~1 -- Behavioural signals.} The customer state agent rebuilds metrics from invoices and payments under a per-customer lock. With payment/due timestamps $\mathrm{paid}_k, \mathrm{due}_k$:
\begin{equation}
\begin{split}
\mathrm{OTR}(c_i) &= \frac{|\{I_k : \mathrm{paid}(I_k)\leq\mathrm{due}(I_k)\}|}{|\mathcal{I}(c_i)|}, \\
\mathrm{AvgDelay}(c_i) &= \frac{1}{|\mathcal{P}|}\sum_{k}\max\bigl((\mathrm{paid}_k - \mathrm{due}_k),\,0\bigr).
\end{split}
\label{eq:behavioural}
\end{equation}

\textbf{Tier~2 -- External signals.} Financial ratios are cached for 90~days as $\mathcal{E}(c_i)\in[0,1]$; litigation scraping scores court exposure over 365~days as $\mathcal{L}(c_i)\in[0,1]$.

\textbf{Tier~3 -- Fusion.} The aggregator fuses external signals:
\begin{equation}
R_{\text{ext}}(c_i) = w_{\text{fin}}\cdot\mathcal{E}(c_i) + w_{\text{lit}}\cdot\mathcal{L}(c_i)\in[0,1],
\label{eq:aggregation}
\end{equation}
with $w_{\text{fin}}{=}0.60$, $w_{\text{lit}}{=}0.40$. Sensitivity is evaluated in Section~\ref{subsec:ablation}.

\subsection{Payment Risk Prediction and Guardrails}\label{subsec:prediction}

The feature vector $X\in\mathbb{R}^{d}$ includes normalised delay, credit utilisation, invoice ratio, late rate, external financial score, litigation index, and rating penalty. A Random Forest ($M{=}50$ trees), pre-calibrated offline, yields $r_{ml}=\frac{1}{M}\sum_{m} h_m(X)$. ACIS-X operates \textbf{label-free at runtime}:
\begin{equation}
r_{final}=\max\bigl(0,\min(1,r_{ml}+r_{policy})\bigr),
\label{eq:prediction}
\end{equation}
where deterministic guardrail $r_{policy}$ enforces rating penalties ($+0.10$~B, $+0.20$~C) and litigation caps ($+0.15$ if $\mathcal{L}{>}0.1$, $+0.25$ if $\mathcal{L}{>}0.5$). Attributions are logged via SHAP TreeExplainer~\cite{lundberg2020local} for auditability.

\subsection{Risk Refinement and Self-Healing}\label{subsec:refinement}

The risk scoring agent forms adjusted risk from prediction $r_{\text{pred}}$:
\begin{equation}
r_{\text{adj}}=\max\Bigl(0,\min\bigl(1,r_{\text{pred}}+0.85\cdot\Delta_{\text{beh}}+0.30\cdot\Delta_{\text{temp}}\bigr)\Bigr),
\label{eq:refinement}
\end{equation}
where $\Delta_{\text{beh}} = 0.40a_{\text{overdue}}{+}0.20a_{\text{pay}}{+}0.15a_{\text{delay}}{+}0.15a_{\text{out}}{+}0.10a_{\text{exp}}$ ranks urgency across normalised factors, and $\Delta_{\text{temp}}{\in}[-0.2,0.3]$ captures trend direction. Agents emit heartbeats every $T_H{=}10$\,s; monitoring detects a 60\,s timeout (6 missed beats), triggering restart from the last committed offset.

\subsection{Model Governance and Stress Harness}\label{subsec:governance}

Layer~6 tracks predictions across sliding windows ($W{=}200$) and computes PSI:
\begin{equation}
\text{PSI} = \sum_{b=1}^{B} (P_b - Q_b) \cdot \ln\!\left(\frac{P_b}{Q_b}\right),
\label{eq:psi}
\end{equation}
where $P_b$, $Q_b$ are baseline/current quantile proportions ($B{=}10$ bins). Values $\text{PSI} \ge 0.10$ trigger drift warnings. A challenger model (linear SGD) is continuously benchmarked against the champion RF.

To evaluate without confidential ledgers, a jump-diffusion stress harness generates scenarios:
\begin{equation}
dZ_t = \kappa (\mu - Z_t)\, dt + \sigma_Z\, dW_t + J_t\, dN_t,
\label{eq:jumpdiff}
\end{equation}
where $Z_t {\in} [0,1]$ is a latent stress factor, $\kappa$ the mean-reversion rate, $\sigma_Z$ volatility, and $J_t\,dN_t$ a Poisson jump (liquidity shocks). Customer delinquency is conditioned on $Z_t$ via Pareto delay mixtures across Calm, Stressed, and Crisis regimes.

%%% ====================================================================
\section{Implementation}\label{sec:implementation}
%%% ====================================================================

ACIS-X is implemented in Python~3.11 (${\sim}55$ files, seven packages). Agents launch as separate OS processes managed by supervisor (max 5 attempts / 2~min). Messages use Pydantic-validated schemas (v1.1) with correlation IDs; validation failures route to a dead-letter queue. Kafka Stream EOS (transactional producers, \texttt{read\_committed} consumers) provides exactly-once for consume-transform-produce loops. External sinks use composite-key upserts (\texttt{ON CONFLICT DO NOTHING}) and LRU deduplication caches.

%%% ====================================================================
\section{Experimental Design and Evaluation}\label{sec:experiments}
%%% ====================================================================

\subsection{Systems Metrics}\label{subsec:sysmetrics}

Tests executed on Windows~11 (i7-12700H, 32\,GB RAM, Docker Kafka~3.6): \textbf{180/180 tests passed} (100\%). Table~\ref{tab:perf} reports system metrics under streaming load across 6 partitions.

\begin{table}[htbp]
\caption{System performance, latency, and scalability.}\label{tab:perf}
\centering\footnotesize
\begin{tabular}{@{}llc@{}}
\toprule
\textbf{Metric} & \textbf{Value} & \textbf{Scope} \\
\midrule
Peak Throughput        & 6{,}420 evt/s             & 6 partitions, batched \\
Latency (P95/P99)      & 16.2\,ms / 28.4\,ms      & Full 5-layer pipeline \\
Consumer Lag           & 0 (steady); $<$15\,ms post-burst & Recovery \\
Dedup (10k repeats)    & 1 action dispatched       & Idempotency \\
Scaling ($N{=}4/6$)    & $3.49{\times}/4.93{\times}$ & Partition-bounded \\
Governance PSI         & 4.1\,ms ($<$0.10 stable)  & Sliding-window \\
Suite Pass Rate        & 180/180 (100\%)           & Full verification \\
\botrule
\end{tabular}
\end{table}

\subsection{Ablation and Sensitivity}\label{subsec:ablation}

Ablation across $n{=}200$ customers over 42 seeds (Table~\ref{tab:ablation}): behavioural refinement provides the dominant lift ($\Delta\rho{=}+0.065$, paired $t{=}37.29$, $p{<}0.0001$, $df{=}41$), while external enrichment contributes marginally ($\Delta\rho{=}+0.005$). In real B2B credit, external financials update quarterly and serve as low-frequency policy guardrails. SHAP adds ${\sim}0.9$\,ms/event.

\begin{table}[htbp]
\caption{Architectural ablation ($n{=}200$, 42 seeds): Spearman $\rho$ with ground truth.}\label{tab:ablation}
\centering\footnotesize
\begin{tabular}{@{}lccc@{}}
\toprule
\textbf{Configuration} & \textbf{$\rho$} & \textbf{$\Delta\rho$} & \textbf{Note} \\
\midrule
Full ACIS-X            & $0.734 \pm 0.035$ & ---      & All layers \\
$-$External enrich.    & $0.728 \pm 0.035$ & $-$0.005 & $p{<}0.0001$ \\
$-$Behav.\ refine.     & $0.669 \pm 0.042$ & $-$0.065 & $p{<}0.0001$ \\
$-$SHAP audit          & (latency)         & $+$0.9\,ms & Overhead \\
\botrule
\end{tabular}
\end{table}

Varying fusion weights ($w_{\text{fin}}/w_{\text{lit}}$: 50/50, 60/40, 70/30) yielded $\rho{=}0.760$, $0.872$, $0.941$ respectively, confirming the pipeline responds predictably to policy weight shifts without online RL.

\subsection{Fault Injection}\label{subsec:fault}

Table~\ref{tab:fault} reports self-healing evaluation. Internal decision latency is sub-millisecond; cooldown prevents restart storms (exactly 1 restart/agent in 5-agent cascades). In distributed deployment, total recovery is governed by the 60\,s heartbeat timeout and Kafka rebalance (3--5\,s), after which consumers resume from committed offsets.

\begin{table}[htbp]
\caption{Fault injection results (single-node decision latency).}\label{tab:fault}
\centering\footnotesize
\begin{tabular}{@{}lccc@{}}
\toprule
\textbf{Scenario} & \textbf{Metric} & \textbf{Value} & \textbf{Assert.} \\
\midrule
Agent crash            & decision lat.  & 0.26\,ms & $<$100\,ms \\
Latency spike (5\,s)   & decision lat.  & 0.10\,ms & $<$100\,ms \\
5-agent cascade        & restarts/agent & 1        & exactly~1 \\
100-cycle endurance    & success rate   & 100\%    & 100/100 \\
100-cycle endurance    & P95 lat.       & 0.23\,ms & $<$50\,ms \\
\botrule
\end{tabular}
\end{table}

\subsection{Predictive Comparison}\label{subsec:predictive}

Table~\ref{tab:ml} benchmarks ACIS-X against supervised learners ($n{=}500$, 42 seeds). Supervised models reach high AUC ($>$0.96) by training on synthetic generator labels unavailable in production. ACIS-X operates label-free, improving upon heuristic rules ($F_1{=}0.757$ vs.\ $0.738$, winning 32/42 seeds) while providing continuous ranking (AUC~$=0.861$) at sub-50\,ms latency.

\begin{table}[htbp]
\caption{Predictive benchmark ($n{=}500$, 42 seeds). Supervised models use ground-truth labels; ACIS-X is label-free.}\label{tab:ml}
\centering\footnotesize
\begin{tabular}{@{}lcc@{}}
\toprule
\textbf{Method} & \textbf{F1} & \textbf{AUC} \\
\midrule
Naive utilisation             & $0.738 \pm 0.038$ & --- \\
ACIS-X streaming (label-free) & $0.757 \pm 0.035$ & $0.861 \pm 0.032$ \\
Log.\ Reg.\ (basic/enriched, sup.) & $0.927/0.965$  & $0.988/0.998$ \\
Random Forest (basic/enriched)      & $0.886/0.895$  & $0.967/0.972$ \\
Gradient Boosting (basic/enriched)  & $0.879/0.888$  & $0.967/0.970$ \\
\botrule
\end{tabular}
\end{table}

%%% ====================================================================
\section{Discussion}\label{sec:discussion}

The evaluation validates ACIS-X as a low-latency, resilient multi-agent architecture for real-time credit intelligence. Behavioural transaction dynamics drive intraday differentiation ($\Delta\rho{=}+0.065$, $p{<}0.0001$), while external ratios act as low-frequency guardrails. PSI and SHAP ($<$1\,ms) run inline without bottlenecking throughput.

\subsection{Limitations}\label{subsec:limitations}

\textbf{Domain Data:} Public retail credit benchmarks lack B2B invoice dynamics; ACIS-X's primary contribution is the situational streaming pipeline evaluated under jump-diffusion tail-risk scenarios. \textbf{Distributed Recovery:} The 0.26\,ms metric reflects single-node decision latency; multi-node recovery is governed by 60\,s heartbeat timeout and Kafka rebalances (3--5\,s). \textbf{Governance:} Automated PSI/challenger signals require coupling with human-in-the-loop credit committee review.

%%% ====================================================================
\section{Conclusion and Future Work}\label{sec:conclusion}
%%% ====================================================================

ACIS-X presents an event-driven multi-agent credit architecture on Kafka featuring three-tier signal fusion, stream-level EOS with idempotent sinks, autonomous model governance, SHAP audit persistence, and heartbeat-driven recovery. Across 42 seeds, the pipeline achieved 6{,}420 events/s peak throughput, P95 16.2\,ms latency, and $4.93{\times}$ partition scaling. Behavioural refinement provides the primary ranking lift ($\Delta\rho{=}+0.065$) with $<$1\,ms SHAP overhead, and fault injection confirms sub-millisecond self-healing decisions. Future work targets field validation via anonymized ERP partnerships and distributed disaster recovery drills.

\backmatter

\bmhead{Acknowledgements}
Faculty and laboratory infrastructure of the Department of CSE, AMC Engineering College, Bengaluru are gratefully acknowledged.

\section*{Declarations}
\small
\noindent\textbf{Funding:} None. \textbf{Competing interests:} The authors declare no competing interests. \textbf{Ethics approval:} Not applicable. \textbf{Data availability:} Synthetic data generated via Section~\ref{subsec:governance}. \textbf{Code availability:} \texttt{https://github.com/nikhilnagendra/ACIS-X}. \textbf{Author contributions:} N.~Nagendra: conceptualization, architecture, implementation, experimentation, manuscript. P.~K~V: supervision. B.~Joshi, J.~K~George, N.~B.~MVS, P.~B.~Acharya: implementation support and review.

\bibliography{references}

\end{document}
